\section{Human-like Agents：能力扩张必须与权限分级同步}

这一部分从“为什么需要 agents”走向“为什么让 agents 使用人类界面”。讲者的答案不是界面更酷，而是现实世界已经为人类建立了浏览器、桌面、表单、登录与支付流程；如果 agent 能在这些界面上操作，就能跨越 API 不完整、服务异构和长尾任务。但直接控制也扩大了动作空间与风险，因此必须把自治级别、确认点和回滚写入产品合同。

\lecturefigure{slide-07-why-agents.jpg}{为什么需要 Agents：单次模型调用不足以承担真实任务}{00:07:10}

\lecturefigure{slide-08-lecture-roadmap.jpg}{课程路线：架构、交互、记忆、通信与未来方向}{00:07:37}

\subsection{从问题清单到交互 thesis}

“Why、How、Ingredients、What can they do” 四个问题给后文提供了教学框架。最终 build 展示讲者的 thesis：人类用自然语言表达意图，AI 操作机器。这里的关键不是自然语言本身，而是把模糊意图变成受约束的动作，并让用户能够观察、修改和撤销。

\lecturefigure{slide-09-building-agents-questions.jpg}{构建 Agent 的四个基本问题}{00:07:54}

\lecturefigure{slide-10-building-agents-thesis.jpg}{交互 thesis：人用自然语言，Agent 操作机器}{00:08:27}

\lecturefigure{slide-11-ai-agents-why-how.jpg}{单次 foundation-model 调用与完整 Agent 系统的区别}{00:08:53}

设任务由 $n$ 个必须成功的步骤组成，第 $i$ 步成功概率为 $p_i$。若错误无法恢复，端到端成功率近似为
\[
P(\text{success})=\prod_{i=1}^{n}p_i.
\]
若每步都是 $p_i=0.98$，20 步任务也只有 $0.98^{20}\approx 0.668$。这解释了为什么聊天中“偶尔犯错”的模型，一旦进入长任务就会显得非常脆弱；agent 系统必须加入验证、重试、分支搜索和恢复，而不是只期待基础模型再提升一点。

\begin{warningbox}{长任务的可靠性不是单步准确率}
平均单步正确率会掩盖关键步骤和风险差异。支付确认、删除文件或发送公开消息的错误成本远高于一次无效搜索；验收应按步骤重要性、错误可逆性和暴露频率分层。
\end{warningbox}

\subsection{为什么追求 human-like interface}

最终完整 slide 给出五个动机：复用为人设计的界面、成为用户的数字延伸、跨越 API 限制、使用简单 click/type 动作，以及从用户交互中持续学习。读图时要同时看到反面：界面为人设计，不等于适合机器；视觉变化、弹窗、广告、登录、反自动化机制和模糊反馈都会制造分布偏移。

\lecturefigure{slide-12-human-like-agent-benefits.jpg}{Human-like agents 的五个动机与隐含风险}{00:11:29}

\teachervoice{00:08:53--00:11:36，讲者强调“一次 LLM call 不够”，并把 human-like agent 定义为能在现有界面上代表用户行动、处理登录与支付、使用 click/type primitives、再从用户处学习。课堂提示：每个优势同时扩大了权限与攻击面。}

\begin{knowledgebox}{术语消化：Agent 能力栈}
\begin{tabularx}{\textwidth}{>{\bfseries}p{0.18\textwidth} X X}
术语 & 解决的问题 & 本讲中的系统关系 \\
\hline
Browser control & 没有稳定 API 时怎样操作服务 & 扩大覆盖面，但引入视觉变化与不可逆动作 \\
Tool calling & 怎样把语言请求映射到结构化调用 & 更可控，但受 schema、权限和错误返回约束 \\
Computer use & 怎样操作桌面、键鼠与多应用 & 动作通用，观测和恢复更困难 \\
Long-term memory & 怎样跨会话保存用户与任务状态 & 支持连续性，同时带来过期、隐私和污染风险 \\
Self-correction & 怎样从失败状态寻找替代路径 & 需要可判定反馈、预算与停止条件 \\
\end{tabularx}
\end{knowledgebox}

\subsection{API 与直接控制是两种不同的风险模型}

API 路线提供结构化 schema、明确错误码和较小动作空间，通常更易审计；直接浏览器或桌面控制覆盖长尾任务，但观察是像素或 accessibility tree，动作含义依赖上下文。二者不是“旧方式与新方式”，而是应按任务风险选择的两种 contract。

读这张对照图时应沿着三个维度判断。第一是\textbf{状态可见性}：API 往往返回结构化字段，直接控制只能从页面推断；第二是\textbf{动作可逆性}：读取接口容易重试，提交表单或点击购买可能产生真实副作用；第三是\textbf{兼容成本}：API 会版本升级，网页也会重排。好的系统不会永久押注一种路线，而会优先使用结构化接口，在缺少覆盖时才降级到浏览器，并对降级动作提高确认等级。

\lecturefigure{slide-13-agent-computer-two-routes.jpg}{Agent--computer interaction 的两条路线：API 与直接控制}{00:12:31}

\begin{importantbox}{动作合同应该显式记录什么}
对每次动作 $a_t$，至少记录目标资源、参数、调用者、权限依据、幂等键、预期后置条件和回滚方法。若动作是网页点击，也应把“点击哪个可访问节点”转成稳定标识，而不是只保存坐标。
\end{importantbox}

\subsection{自治等级：能做与应该做不是同一个问题}

讲者借用自动驾驶等级图来讨论 agents。更高自治不是简单减少用户点击，而是把环境监控、异常判断和 fallback 责任从人移向系统。每提升一级，都要说明谁监控、谁确认、谁负责终止，以及系统失效时是否仍有安全状态。

\lecturefigure{slide-14-five-levels-autonomy.jpg}{从驾驶自动化类比 Agent 自治等级}{00:13:50}

\begin{warningbox}{自动驾驶类比的边界}
网页任务与道路驾驶共享“感知--决策--行动--接管”结构，但风险、时延和法律责任不同。类比可帮助定义 fallback，不应把某个 agent 产品直接宣传为“L4/L5”。
\end{warningbox}

\subsection{本章小结}

Human-like agents 的价值是复用现实接口并代表用户完成长尾流程；它们的风险也来自同一来源。API 与直接控制应按可控性、覆盖面和可恢复性选择，自治等级必须与监控、确认和接管责任同步设计。下一章用 REAL Bench 说明如何把这些复杂性放进可复现评测。

